\documentclass[a4paper]{article}     
\usepackage{amsfonts}
\usepackage{amsmath}
\usepackage{amssymb}
\usepackage{graphicx}
\setlength{\parindent}{0pt}

\newcommand{\fx}{\dfrac{\sqrt{x+7}-3}{x-2}}
\newcommand{\llim}[1]{ \lim\limits_{x \rightarrow #1} }

\begin{document}

\noindent \large Auteur: Milo Rampelbergh \\
\noindent \large Studierichting: Informatica\\
\noindent \large Oefeningenreeks 2D nr. 11.1\\

\medskip

\normalsize

%Begin oefening

$ f(x) = \fx $ \\

$ \llim{2} \fx = \dfrac{\sqrt{9}-3}{0} = \dfrac{0}{0} = \infty $ \\

$ = \llim{2} \dfrac{(\sqrt{x+7}-3)(\sqrt{x+7}+3)}{(x-2)(\sqrt{x+7}+3)} $ \\

$ = \llim{2} \dfrac{x+7-9}{(x-2)(\sqrt{x+7}+3)} $ \\

$ = \llim{2} \dfrac{1}{\sqrt{x+7}+3} = \dfrac{1}{6} $ \\

%Einde oefening

\begin{thebibliography}{999}
%\bibitem{a} Referentie
\end{thebibliography}
\end{document} 